% Source: https://artofproblemsolving.com/community/c3h2013830_real_solution_to_this_problem
% Stable identity: aops-topic-2013830-q01
Reflect $K$ in $M$ to get $K'$, so $M$ is the midpoint of both $BC$ and $KK'$. Then $BKCK'$ is a parallelogram (its diagonals bisect each other), and opposite angles of a parallelogram are equal:
\[
\angle BK'C=\angle BKC=180^\circ-\angle BAC .
\]

\begin{asy}
usepackage("amsmath");
size(5.5cm);
defaultpen(fontsize(10pt));
pair A=(1.1,3.6), B=(0,0), C=(4.4,0);
pair M=(B+C)/2;
pair mAB=(A+B)/2, mAC=(A+C)/2;
pair O=extension(mAB,mAB+rotate(90)*(B-A),mAC,mAC+rotate(90)*(C-A));
real R=abs(O-A);
real t=abs(B-M)*abs(C-M)/abs(A-M);
pair Kp=M+t*unit(M-A), K=M+t*unit(A-M);
draw(circle(O,R),gray(0.5));
draw(A--B--C--cycle,linewidth(0.9));
draw(A--Kp,blue);
draw(B--K--C,red); draw(B--Kp--C,red+dashed);
dot(A); dot(B); dot(C); dot(M); dot(K); dot(Kp);
label("$A$",A,N); label("$B$",B,W); label("$C$",C,E); label("$M$",M,SE);
label("$K$",K,W); label("$K'$",Kp,S);
label("\mbox{}",(min(currentpicture,user=true)+max(currentpicture,user=true))/2,black); // keeps the text after the figure black when LaTeX typesets the labels inline
\end{asy}

Since $K$ lies on ray $MA$, its reflection $K'$ lies on the ray from $M$ opposite to $A$, so $A$ and $K'$ are on opposite sides of line $BC$. In quadrilateral $ABK'C$ the opposite angles at $A$ and $K'$ add to $180^\circ$, so $A$, $B$, $K'$, $C$ lie on a circle.

The chords $AK'$ and $BC$ of this circle meet at $M$, so by the power of the point $M$,
\[
AM\cdot MK'=BM\cdot MC=\left(\dfrac{\sqrt7}{2}\right)^2=\dfrac{7}{4} .
\]
As $MK'=MK$,
\[
AM\cdot MK=\boxed{\dfrac{7}{4}}.
\]
(The poster's guess with $\angle BAC=90^\circ$ gives the same value, because the product does not depend on the shape of the triangle; in that special case $K=A$.)
